\section{기대값 기반의 결정 이론}

기대값 기반의 결정 이론은 어떤 선택의 기대값이 높을수록 그 선택을 하는 것이 합리적이라는 이론이다. (조사 필요)
어떤 선택의 기대값이 0 이상이면 그 선택을 하는 것이 합리적이라고 할 수 있다.

\section{St. Petersburg Paradox}

St. Petersburg Paradox는 기대값 계산 방식이 상식에 어긋나는 결과를 도출함을 보여주는 문제이다. St. Petersburg Game은 다음과 같은 규칙을 가진 게임이다. 참가자는 동전을 던져서 앞면이 나올 때까지 반복하고, 첫 번째로 앞면이 나오기까지 동전을 던진 횟수에 따라 상금을 받는다. 이때 동전을 던진 횟수를 X라고 하면 상금은 $2^X$원이다.

St. Petersburg Game에서 참가자가 얻을 수 있는 상금의 기대값은 다음과 같다.
\[
  \mathbf{E}[X] = \sum_{x=1}^{\infty} \frac{1}{2^x} \cdot 2^x = \sum_{x=1}^{\infty} 1 = \infty
\]

이론적으로 St. Petersburg Game의 기대값은 무한대이므로, 참가비가 아무리 높아도 그 금액이 유한하다면 이 게임에 참여하는 것이 합리적이어야 한다. 하지만 현실적으로는 사람들은 이 게임에 큰 돈을 지불하지 않으려고 한다.

\section{St. Petersburg Paradox에 대한 기존 해석}
